% ============================================================================
%  PRESENTACIÓN PEOPLE ANALYTICS – FINANCORP CHILE
%  Estilo: Beamer moderno con diseño ejecutivo premium
% ============================================================================
\documentclass[aspectratio=169,12pt]{beamer}

% ── Fuentes ──────────────────────────────────────────────────────────────────
\usepackage[T1]{fontenc}
\usepackage[utf8]{inputenc}
\usepackage{lmodern}
\usepackage{microtype}

% ── Idioma ───────────────────────────────────────────────────────────────────
\usepackage[spanish]{babel}

% ── Matemáticas y símbolos ───────────────────────────────────────────────────
\usepackage{amsmath,amssymb}

% ── Gráficos y colores ───────────────────────────────────────────────────────
\usepackage{graphicx}
\usepackage{tikz}
\usetikzlibrary{shapes.geometric,arrows.meta,positioning,calc,decorations.pathreplacing}
\usepackage{pgfplots}
\pgfplotsset{compat=1.18}

% ── Tablas ───────────────────────────────────────────────────────────────────
\usepackage{booktabs}
\usepackage{array}
\usepackage{multirow}
\usepackage{colortbl}

% ── Íconos y extras ──────────────────────────────────────────────────────────
\usepackage{fontawesome5}
\usepackage{xcolor}
\usepackage{tcolorbox}
\tcbuselibrary{skins,breakable,raster}
\usepackage{calc}

% ============================================================================
%  PALETA DE COLORES CORPORATIVA
% ============================================================================
\definecolor{NavyDeep}    {HTML}{0D1B2A}
\definecolor{NavyMid}     {HTML}{1B3A5C}
\definecolor{NavyLight}   {HTML}{2E6DA4}
\definecolor{Cyan}        {HTML}{00B4D8}
\definecolor{Teal}        {HTML}{0077B6}
\definecolor{Gold}        {HTML}{F4A261}
\definecolor{Amber}       {HTML}{E76F51}
\definecolor{GreenOK}     {HTML}{2EC4B6}
\definecolor{RedRisk}     {HTML}{E63946}
\definecolor{GrayLight}   {HTML}{F0F4F8}
\definecolor{GrayMid}     {HTML}{8899AA}
\definecolor{White}       {HTML}{FFFFFF}
\definecolor{TextDark}    {HTML}{1A1A2E}

% ============================================================================
%  TEMA BASE Y PERSONALIZACIÓN
% ============================================================================
\usetheme{Madrid}
\usecolortheme[named=NavyDeep]{structure}

\setbeamercolor{palette primary}    {bg=NavyDeep,  fg=White}
\setbeamercolor{palette secondary}  {bg=NavyMid,   fg=White}
\setbeamercolor{palette tertiary}   {bg=NavyLight, fg=White}
\setbeamercolor{palette quaternary} {bg=Cyan,      fg=NavyDeep}
\setbeamercolor{title}              {fg=White}
\setbeamercolor{subtitle}           {fg=Cyan}
\setbeamercolor{frametitle}         {bg=NavyDeep,fg=White}
\setbeamercolor{block title}        {bg=NavyLight,fg=White}
\setbeamercolor{block body}         {bg=GrayLight,fg=TextDark}
\setbeamercolor{block title alerted}{bg=Amber,fg=White}
\setbeamercolor{block body alerted} {bg=GrayLight,fg=TextDark}
\setbeamercolor{block title example}{bg=GreenOK,fg=White}
\setbeamercolor{block body example} {bg=GrayLight,fg=TextDark}
\setbeamercolor{normal text}        {fg=TextDark}
\setbeamercolor{itemize item}       {fg=Cyan}
\setbeamercolor{itemize subitem}    {fg=NavyLight}
\setbeamercolor{footline}           {bg=NavyDeep,fg=GrayMid}

\setbeamerfont{title}{size=\LARGE,series=\bfseries}
\setbeamerfont{subtitle}{size=\large}
\setbeamerfont{frametitle}{size=\large,series=\bfseries}
\setbeamerfont{section in toc}{series=\bfseries}

\setbeamertemplate{itemize item}{\faChevronRight}
\setbeamertemplate{itemize subitem}{\faAngleRight}
\setbeamertemplate{navigation symbols}{}

% Barra de progreso inferior
\makeatletter
\addtobeamertemplate{footline}{}{%
  \begin{tikzpicture}[overlay,remember picture]
    \fill[Cyan,opacity=0.85]
      (current page.south west)
      rectangle
      ($(current page.south west)+(\insertframenumber/\inserttotalframenumber*\paperwidth,2pt)$);
  \end{tikzpicture}%
}
\makeatother

% ============================================================================
%  COMANDOS AUXILIARES
% ============================================================================

% Caja de KPI  \kpibox{color}{icono}{valor}{etiqueta}
\newcommand{\kpibox}[4]{%
  \begin{tcolorbox}[
    enhanced,
    colback=#1!10!White,
    colframe=#1,
    arc=4pt,
    boxrule=1.5pt,
    left=4pt, right=4pt, top=3pt, bottom=3pt,
    width=\linewidth,
    halign=center
  ]
    \centering
    {\color{#1}\Large #2}\\[2pt]
    {\color{NavyDeep}\Large\bfseries #3}\\[1pt]
    {\color{GrayMid}\footnotesize #4}
  \end{tcolorbox}%
}

% Caja de hallazgo  \hallabox{color}{titulo}{texto}
\newcommand{\hallabox}[3]{%
  \begin{tcolorbox}[
    enhanced,
    colback=White,
    colframe=#1,
    arc=3pt,
    boxrule=1pt,
    borderline west={3pt}{0pt}{#1},
    left=6pt, right=4pt, top=3pt, bottom=3pt
  ]
    {\color{#1}\small\bfseries #2}\\[-2pt]
    {\small #3}
  \end{tcolorbox}%
}

% ============================================================================
%  META-DATOS
% ============================================================================
\title[People Analytics · FinanCorp Chile]{%
  \textbf{People Analytics}\\
  {\large Diagnóstico Predictivo de Rotación y Fuga de Talento}}
\subtitle{FinanCorp Chile · Análisis Estratégico Mayo 2026}
\date{Mayo 2026}
\author{Gerencia de Personas · Área de Analytics}
\institute{FinanCorp Chile}

% ============================================================================
%  DOCUMENTO
% ============================================================================
\begin{document}

% ─────────────────────────────────────────────────────────────────────────────
%  PORTADA
% ─────────────────────────────────────────────────────────────────────────────
{
\setbeamertemplate{headline}{}
\setbeamertemplate{footline}{}
\begin{frame}[plain]
  \begin{tikzpicture}[overlay,remember picture]
    % Fondo degradado
    \fill[NavyDeep] (current page.south west) rectangle (current page.north east);
    \fill[NavyMid,opacity=0.5]
      ([xshift=12cm]current page.south west) --
      (current page.north east) --
      ([yshift=-4cm]current page.north east) -- cycle;
    % Acento superior izquierdo
    \fill[Cyan,opacity=0.7]
      (current page.north west) rectangle
      ([xshift=4pt,yshift=-3pt]current page.north west |- {(0,0)});
    \draw[Cyan,line width=2pt]
      ([yshift=-0.06\paperheight]current page.north west) --
      ([yshift=-0.06\paperheight]current page.north east);
    \draw[Teal,line width=0.5pt,opacity=0.4]
      ([yshift=-0.065\paperheight]current page.north west) --
      ([yshift=-0.065\paperheight]current page.north east);

    % Patrón de puntos decorativo
    \foreach \x in {0.7,0.75,...,0.98}
      \foreach \y in {0.05,0.10,...,0.55}
        \fill[White,opacity=0.04] ([xshift=\x\paperwidth,yshift=\y\paperheight]current page.south west) circle(1pt);

    % Línea acento inferior
    \draw[Gold,line width=2pt]
      ([yshift=0.12\paperheight]current page.south west) --
      ([yshift=0.12\paperheight,xshift=0.45\paperwidth]current page.south west);
  \end{tikzpicture}

  \vspace{0.8cm}
  \begin{columns}[T]
    \column{0.62\textwidth}
      {\color{Cyan}\footnotesize\bfseries\MakeUppercase{FinanCorp Chile · Mayo 2026}}
      \vspace{0.2cm}

      {\color{White}\fontsize{26}{30}\selectfont\bfseries People\\Analytics}
      \vspace{0.15cm}

      {\color{GrayLight}\fontsize{13}{16}\selectfont
        Diagnóstico Predictivo de Rotación,\\
        Retención y Fuga de Talento}
      \vspace{0.35cm}

      \begin{tikzpicture}
        \draw[Cyan,line width=1.5pt] (0,0) -- (6,0);
        \draw[Gold,line width=0.5pt] (0,-3pt) -- (3,-3pt);
      \end{tikzpicture}
      \vspace{0.3cm}

      {\color{GrayMid}\small
        \faUsers\ \ 3.200 colaboradores \quad
        \faCalendar\ \ Mayo 2026}

    \column{0.36\textwidth}
      \vspace{0.4cm}
      % Panel de KPIs en portada
      \begin{tcolorbox}[
        enhanced, colback=NavyMid, colframe=Cyan,
        arc=6pt, boxrule=1pt,
        left=6pt, right=6pt, top=6pt, bottom=6pt
      ]
        \centering
        {\color{Cyan}\footnotesize\bfseries MÉTRICAS CLAVE}
        \vspace{4pt}

        \begin{tabular}{@{}l@{\hspace{4pt}}r@{}}
          {\color{GrayMid}\scriptsize Rotación actual} &
            {\color{RedRisk}\bfseries 22\%} \\[3pt]
          {\color{GrayMid}\scriptsize Salidas/año} &
            {\color{White}\bfseries 703} \\[3pt]
          {\color{GrayMid}\scriptsize Costo anual} &
            {\color{Gold}\bfseries \$2.1M USD} \\[3pt]
          {\color{GrayMid}\scriptsize Modelo ROC-AUC} &
            {\color{GreenOK}\bfseries 67.7\%} \\[3pt]
          {\color{GrayMid}\scriptsize Ahorro potencial} &
            {\color{GreenOK}\bfseries \$1.3M USD} \\
        \end{tabular}
      \end{tcolorbox}
  \end{columns}
\end{frame}
}

% ─────────────────────────────────────────────────────────────────────────────
%  ÍNDICE
% ─────────────────────────────────────────────────────────────────────────────
\begin{frame}{Agenda}
  \begin{columns}[T]
    \column{0.48\textwidth}
      \tableofcontents[sections={1-3}]
    \column{0.48\textwidth}
      \tableofcontents[sections={4-5}]
  \end{columns}
\end{frame}

% ============================================================================
\section[Contexto]{Contexto y Problemática}
% ============================================================================

\begin{frame}[shrink=5]{El Desafío de FinanCorp Chile}
  \begin{columns}[T]
    \column{0.52\textwidth}
      \textbf{\color{NavyDeep}Situación Actual}
      \vspace{3pt}
      \begin{itemize}\setlength\itemsep{2pt}
        \item Empresa de servicios financieros con \textbf{3.200 colaboradores}
        \item Rotación del \textbf{22\%} anual vs.\ benchmark industria \textbf{15--18\%}
        \item \textbf{703 salidas anuales} con costo promedio \$3.000 USD/persona
        \item Concentración crítica en Tecnología, Riesgo y Operaciones
      \end{itemize}
      \vspace{4pt}
      \hallabox{RedRisk}{\faExclamationTriangle\ \ Brecha vs.\ Industria}{%
        Rotación supera en \textbf{4--7 pp} el benchmark del sector financiero,
        representando un riesgo operacional y financiero significativo.}

    \column{0.44\textwidth}
      \begin{tcolorbox}[
        enhanced, colback=GrayLight, colframe=NavyLight,
        arc=5pt, boxrule=1pt, title={\color{White}\bfseries Impacto Financiero},
        attach boxed title to top left={yshift=-2mm,xshift=5mm},
        boxed title style={colback=NavyLight}
      ]
        \centering
        \begin{tikzpicture}
          \begin{axis}[
            ybar, width=4.8cm, height=3.8cm,
            bar width=13pt,
            xtick={1,2,3},
            xticklabels={A\~no 1, A\~no 2, A\~no 3},
            ymin=0, ymax=2500,
            ytick={500,1000,1500,2000,2500},
            yticklabels={\$0.5M,\$1M,\$1.5M,\$2M,\$2.5M},
            tick label style={font=\tiny},
            axis lines=left,
            ymajorgrids=true, grid style={dotted,gray!30},
            every axis plot/.append style={fill=RedRisk!70,draw=none}
          ]
            \addplot coordinates {(1,2109) (2,2109) (3,2109)};
          \end{axis}
        \end{tikzpicture}\\
        {\scriptsize\color{GrayMid}Costo anual sin intervención}
      \end{tcolorbox}
  \end{columns}
\end{frame}

% ─────────────────────────────────────────────────────────────────────────────
\begin{frame}{Objetivos y Metodología}
  \begin{columns}[T]
    \column{0.48\textwidth}
      \begin{block}{\faTarget\ \ Objetivos}
        \begin{enumerate}\setlength\itemsep{1pt}
          \item Identificar factores \textbf{causales} de rotación
          \item Construir \textbf{modelo predictivo} (ML)
          \item \textbf{Segmentar} colaboradores por nivel de riesgo
          \item Priorizar \textbf{intervenciones} con ROI
          \item Proyectar \textbf{impacto financiero} esperado
        \end{enumerate}
      \end{block}

    \column{0.48\textwidth}
      \begin{block}{\faCogs\ \ Metodología}
        \begin{itemize}\setlength\itemsep{1pt}
          \item \textbf{EDA:} Estadísticas descriptivas y correlaciones
          \item \textbf{Scoring:} Índice de riesgo ponderado (0--100)
          \item \textbf{ML:} Regresión Logística vs.\ Random Forest
          \item \textbf{Priorización:} Matriz impacto × urgencia
          \item \textbf{ROI:} Análisis costo-beneficio a 3 años
        \end{itemize}
      \end{block}
  \end{columns}
  \vspace{2pt}
  \begin{tcolorbox}[
    enhanced, colback=Cyan!8, colframe=Cyan,
    arc=3pt, boxrule=0.8pt, left=4pt, right=4pt, top=2pt, bottom=2pt
  ]
    {\color{NavyDeep}\small\bfseries Dataset:} {\small 3.200 colaboradores · 9 variables predictoras · variable objetivo: \texttt{Rotación}}
  \end{tcolorbox}
\end{frame}

% ============================================================================
\section[EDA]{Análisis Exploratorio}
% ============================================================================

\begin{frame}[shrink=3]{Dataset: Estadísticas Descriptivas}
  \begin{columns}[T]
    \column{0.52\textwidth}
      \small
      \begin{tabular}{@{}lrr@{}}
        \toprule
        \rowcolor{NavyDeep}\color{White}\textbf{Variable} &
          \color{White}\textbf{Media} &
          \color{White}\textbf{Descripción} \\
        \midrule
        \rowcolor{GrayLight}Antigüedad (años) & 3.00 & Dist.\ exponencial \\
        Desempeño (1--5)  & 3.17 & Escala evaluación \\
        \rowcolor{GrayLight}Clima (1--5)      & 3.06 & Ambiente laboral \\
        Compromiso (1--5) & 2.20 & \textbf{Factor crítico} \\
        \rowcolor{GrayLight}Ausentismo (días) & 4.89 & Indicador estrés \\
        Capacitación (cur/año) & 2.1 & Desarrollo \\
        \rowcolor{GrayLight}\textbf{Rotación} & \textbf{22\%} & \textbf{Variable objetivo} \\
        \bottomrule
      \end{tabular}

    \column{0.44\textwidth}
      \centering
      {\small\bfseries\color{NavyDeep}Distribución por Área}
      \vspace{2pt}
      \begin{tikzpicture}
        \begin{axis}[
          xbar, width=5cm, height=4.5cm,
          bar width=7pt,
          symbolic y coords={Administración,Legal,Personas,Finanzas,Riesgo,Clientes,Operaciones,Tecnología},
          ytick=data,
          xmin=0, xmax=900,
          tick label style={font=\tiny},
          axis lines=left,
          xmajorgrids=true, grid style={dotted,gray!30},
          every axis plot/.append style={fill=NavyLight!80,draw=none}
        ]
          \addplot coordinates {
            (32,Administración)(64,Legal)(128,Personas)
            (256,Finanzas)(576,Riesgo)(640,Clientes)
            (704,Operaciones)(800,Tecnología)
          };
        \end{axis}
      \end{tikzpicture}
  \end{columns}
\end{frame}

% ─────────────────────────────────────────────────────────────────────────────
\begin{frame}[shrink=5]{Correlaciones con Rotación}
  \begin{columns}[T]
    \column{0.50\textwidth}
      \centering
      \begin{tikzpicture}
        \begin{axis}[
          xbar, width=6cm, height=5cm,
          bar width=8pt,
          symbolic y coords={Antigüedad,Movilidad,Capacitación,Desempeño,Clima,Compromiso,Ausentismo,Área Crít.},
          ytick={Antigüedad,Movilidad,Capacitación,Desempeño,Clima,Compromiso,Ausentismo,Área Crít.},
          xmin=-0.2, xmax=0.25,
          tick label style={font=\scriptsize},
          axis lines=left,
          xmajorgrids=true, grid style={dotted,gray!25},
          xtick={-0.2,-0.1,0,0.1,0.2},
          xticklabels={-0.2,-0.1,0,0.1,0.2},
          title={\scriptsize Correlación de Pearson con Rotación},
          title style={font=\scriptsize,color=NavyMid}
        ]
          % Negativas (protectoras) en azul
          \addplot[fill=NavyLight!70,draw=none, bar shift=0pt] coordinates {
            (-0.002,Antigüedad)(-0.009,Movilidad)(-0.017,Capacitación)
            (-0.063,Desempeño)(-0.067,Clima)(-0.099,Compromiso)
          };
          % Positiva de riesgo
          \addplot[fill=RedRisk!70,draw=none, bar shift=0pt] coordinates {
            (0.114,Ausentismo)
          };
          % Más fuerte positivo
          \addplot[fill=Amber!90,draw=none, bar shift=0pt] coordinates {
            (0.185,Área Crít.)
          };
        \end{axis}
      \end{tikzpicture}

    \column{0.46\textwidth}
      \vspace{2pt}
      \hallabox{Amber}{\faExclamationCircle\ \ Factor N°1: Área Crítica}{%
        Correlación \textbf{+0.185}: pertenecer a un
        área crítica es el predictor más fuerte
        de rotación en el modelo.}
      \vspace{3pt}
      \hallabox{RedRisk}{\faChartLine\ \ Factor N°2: Ausentismo}{%
        Correlación \textbf{+0.114}. Señal de alerta
        temprana de intención de abandono.}
      \vspace{3pt}
      \hallabox{NavyLight}{\faHeart\ \ Factor N°3: Compromiso}{%
        Correlación \textbf{--0.099}. Bajo compromiso
        asociado a mayor probabilidad de rotar.}
  \end{columns}
\end{frame}

% ─────────────────────────────────────────────────────────────────────────────
\begin{frame}[shrink=12]{Panorama Global de Rotación}
  \begin{columns}[T]
    \column{0.3\textwidth}
      \kpibox{RedRisk}{\faArrowUp}{22\%}{Rotación Global}
      \vspace{3pt}
      \kpibox{Amber}{\faUsers}{703}{Salidas / Año}
      \vspace{3pt}
      \kpibox{GreenOK}{\faDollarSign}{\$2.1M}{Costo USD / Año}

    \column{0.67\textwidth}
      \small
      \begin{tabular}{@{}lrrrr@{}}
        \toprule
        \rowcolor{NavyDeep}
          \color{White}\textbf{Categoría} &
          \color{White}\textbf{N} &
          \color{White}\textbf{Rotación} &
          \color{White}\textbf{Salidas} &
          \color{White}\textbf{USD} \\
        \midrule
        \rowcolor{RedRisk!15}\textbf{Áreas Críticas}
          & 2.720 & \textbf{25.2\%} & 685 & \$2.06M \\
        Otras Áreas
          & 480 & 3.8\% & 18 & \$54K \\
        \midrule
        \rowcolor{GrayLight}\textbf{TOTAL}
          & \textbf{3.200} & \textbf{22.0\%} & \textbf{703} & \textbf{\$2.11M} \\
        \bottomrule
      \end{tabular}
      \vspace{4pt}

      \begin{tcolorbox}[
        enhanced, colback=RedRisk!8, colframe=RedRisk,
        arc=3pt, boxrule=0.8pt, left=4pt, right=4pt, top=2pt, bottom=2pt
      ]
        {\color{RedRisk}\small\bfseries \faExclamationTriangle\ \ Diferencial Crítico:}
        {\small La rotación en áreas críticas es \textbf{6.6× superior}
        (25.2\% vs.\ 3.8\%). Diferencia: \textbf{21.4 pp}.}
      \end{tcolorbox}
  \end{columns}
\end{frame}

% ─────────────────────────────────────────────────────────────────────────────
\begin{frame}{Rotación por Antigüedad y Desempeño}
  \begin{columns}[T]
    \column{0.48\textwidth}
      {\small\bfseries\color{NavyDeep}Por Tramo de Antigüedad}
      \vspace{2pt}
      \begin{tikzpicture}
        \begin{axis}[
          ybar, width=5.5cm, height=3.8cm,
          bar width=12pt,
          xtick={1,2,3,4,5},
          xticklabels={0--1a,1--2a,2--5a,5--10a,+10a},
          ymin=0, ymax=40,
          ytick={0,10,20,30,40},
          yticklabels={0\%,10\%,20\%,30\%,40\%},
          tick label style={font=\tiny},
          axis lines=left,
          ymajorgrids=true, grid style={dotted,gray!30}
        ]
          \addplot[fill=RedRisk!80,draw=none]
            coordinates {(5,24.1)};
          \addplot[fill=Gold!80,draw=none]
            coordinates {(1,22.2)(2,21.5)(3,22.7)};
          \addplot[fill=GreenOK!80,draw=none]
            coordinates {(4,20.3)};
        \end{axis}
      \end{tikzpicture}\\
      {\tiny\color{GrayMid}Rotación uniforme entre tramos (20--24\%)}

    \column{0.48\textwidth}
      {\small\bfseries\color{NavyDeep}Por Nivel de Desempeño}
      \vspace{2pt}
      \begin{tikzpicture}
        \begin{axis}[
          ybar, width=5.5cm, height=3.8cm,
          bar width=14pt,
          xtick={1,2,3,4},
          xticklabels={Bajo,Medio,Bueno,Excelente},
          ymin=0, ymax=40,
          ytick={0,10,20,30,40},
          yticklabels={0\%,10\%,20\%,30\%,40\%},
          tick label style={font=\tiny},
          axis lines=left,
          ymajorgrids=true, grid style={dotted,gray!30}
        ]
          \addplot[fill=RedRisk!80,draw=none]
            coordinates {(1,26.3)};
          \addplot[fill=Gold!80,draw=none]
            coordinates {(2,22.5)(3,22.4)};
          \addplot[fill=GreenOK!80,draw=none]
            coordinates {(4,17.1)};
        \end{axis}
      \end{tikzpicture}\\
      {\tiny\color{GrayMid}Alto desempeño retiene (17\%); bajo desempeño rota más (26\%)}
  \end{columns}
  \vspace{2pt}
  \hallabox{Gold}{\faLightbulb\ \ Hallazgo Clave}{%
    Bajo desempeño: 26.3\% rotación. Excelente desempeño: 17.1\% (menor rotación, mejores condiciones).}
\end{frame}

% ─────────────────────────────────────────────────────────────────────────────
\begin{frame}{Desglose por Área Funcional}
  \begin{columns}[T]
    \column{0.56\textwidth}
      \setlength{\tabcolsep}{4pt}%
      \resizebox{\linewidth}{!}{%
      \begin{tabular}{@{}lrrrrr@{}}
        \toprule
        \rowcolor{NavyDeep}
          \color{White}\textbf{Área} &
          \color{White}\textbf{N} &
          \color{White}\textbf{Rota.} &
          \color{White}\textbf{Salidas} &
          \color{White}\textbf{Costo} &
          \color{White}\textbf{\%~Cía.} \\
        \midrule
        \rowcolor{RedRisk!20}Atención Cli.  & 640  & 27.8\% & 178 & \$534K & 25.3\% \\
        \rowcolor{RedRisk!12}Tecnología     & 800  & 25.5\% & 204 & \$612K & 29.0\% \\
        \rowcolor{RedRisk!8} Riesgo/Cump.  & 576  & 24.0\% & 138 & \$414K & 19.6\% \\
        \rowcolor{RedRisk!5} Operaciones   & 704  & 23.4\% & 165 & \$495K & 23.5\% \\
        \rowcolor{GreenOK!10}Personas       & 128  &  4.7\% &   6 & \$18K  &  0.9\% \\
        \rowcolor{GreenOK!10}Finanzas       & 256  &  4.3\% &  11 & \$33K  &  1.6\% \\
        \rowcolor{GreenOK!10}Legal          &  64  &  1.6\% &   1 & \$3K   &  0.1\% \\
        \rowcolor{GreenOK!10}Administración &  32  &  0.0\% &   0 & \$0K   &  0.0\% \\
        \midrule
        \rowcolor{GrayLight}\textbf{TOTAL} & \textbf{3.200} & \textbf{22.0\%}
          & \textbf{703} & \textbf{\$2.1M} & \textbf{100\%} \\
        \bottomrule
      \end{tabular}%
      }

    \column{0.40\textwidth}
      {\scriptsize\bfseries\color{NavyDeep}Tasa de Rotación por Área (\%)}
      \vspace{5pt}
      \begin{tikzpicture}
        % Líneas de grilla verticales
        \draw[gray!22,thin] (0.857cm,-0.08cm)--(0.857cm,3.06cm);
        \draw[gray!22,thin] (1.714cm,-0.08cm)--(1.714cm,3.06cm);
        \draw[gray!22,thin] (2.571cm,-0.08cm)--(2.571cm,3.06cm);
        % Barras (de abajo: Legal → Tecnología)
        % Escala: 1% = 3cm/35 = 0.0857cm
        \fill[GreenOK!80]  (0,0.00cm) rectangle (0.14cm, 0.25cm); % Legal  1.6%
        \fill[GreenOK!80]  (0,0.52cm) rectangle (0.37cm, 0.77cm); % Finanzas 4.3%
        \fill[Gold!65]     (0,1.04cm) rectangle (2.00cm, 1.29cm); % Operaciones 23.4%
        \fill[Gold!85]     (0,1.56cm) rectangle (2.06cm, 1.81cm); % Riesgo 24.0%
        \fill[RedRisk!65]  (0,2.08cm) rectangle (2.19cm, 2.33cm); % Tecnología 25.5%
        \fill[RedRisk!85]  (0,2.60cm) rectangle (2.38cm, 2.85cm); % Atención 27.8%
        % Etiquetas eje Y
        \node[anchor=east,font=\tiny] at (-0.07cm,0.125cm) {Legal};
        \node[anchor=east,font=\tiny] at (-0.07cm,0.645cm) {Finanzas};
        \node[anchor=east,font=\tiny] at (-0.07cm,1.165cm) {Operaciones};
        \node[anchor=east,font=\tiny] at (-0.07cm,1.685cm) {Riesgo};
        \node[anchor=east,font=\tiny] at (-0.07cm,2.205cm) {Tecnología};
        \node[anchor=east,font=\tiny] at (-0.07cm,2.725cm) {Atención};
        % Eje Y
        \draw[gray!50,thin] (0,-0.05cm)--(0,2.97cm);
        % Eje X con flecha
        \draw[gray!50,thin,->] (-0.03cm,-0.05cm)--(3.1cm,-0.05cm);
        % Ticks y etiquetas eje X
        \foreach \xv/\lbl in {0/0, 0.857/10, 1.714/20, 2.571/30} {
          \draw[gray!50,thin] (\xv cm,-0.05cm)--(\xv cm,-0.13cm);
          \node[font=\tiny,anchor=north] at (\xv cm,-0.15cm) {\lbl};
        }
      \end{tikzpicture}
  \end{columns}
\end{frame}

% ─────────────────────────────────────────────────────────────────────────────
\begin{frame}[shrink=10]{Modelo de Score de Riesgo (0--100)}
  \begin{columns}[T]
    \column{0.50\textwidth}
      {\footnotesize La puntuación de riesgo es una combinación ponderada de los
      factores predictores más relevantes:}
      \vspace{1pt}

      \[
      \textstyle\text{Score} = \sum_{i=1}^{n} w_i \cdot f_i(x_i)
      \]

      \vspace{1pt}
      \footnotesize
      \begin{tabular}{@{}lrr@{}}
        \toprule
        \rowcolor{NavyDeep}\color{White}\textbf{Factor} &
          \color{White}\textbf{Peso} &
          \color{White}\textbf{Umbral} \\
        \midrule
        \rowcolor{RedRisk!15}Compromiso Bajo & 35\% & $<2.5/5$ \\
        \rowcolor{Amber!15}  Ausentismo Alto & 25\% & $>12$ días/año \\
        \rowcolor{Gold!15}   Bajo Desempeño  & 20\% & $<2.5/5$ \\
        Falta Capacitación   & 10\% & $<1$ curso/año \\
        \rowcolor{GrayLight} Estancamiento   &  5\% & Sin movilidad 3+a \\
        Área Crítica         &  5\% & Binaria \\
        \bottomrule
      \end{tabular}

    \column{0.46\textwidth}
      \vspace{1pt}
      \footnotesize
      \begin{tabular}{@{}lrrr@{}}
        \toprule
        \rowcolor{NavyDeep}
          \color{White}\textbf{Nivel} &
          \color{White}\textbf{N} &
          \color{White}\textbf{\%} &
          \color{White}\textbf{Rot.} \\
        \midrule
        \rowcolor{GreenOK!15}  Bajo (0--25)   &   724 & 22.6\% & 16.9\% \\
        \rowcolor{Gold!15}     Medio (25--50) & 1.385 & 43.3\% & 21.1\% \\
        \rowcolor{Amber!20}    Alto (50--75)  &   966 & 30.2\% & 27.1\% \\
        \rowcolor{RedRisk!20}  Muy Alto (75+) &    79 &  2.5\% & 30.4\% \\
        \midrule
        \textbf{TOTAL} & \textbf{3.200} & \textbf{100\%} & \textbf{22.0\%} \\
        \bottomrule
      \end{tabular}
      \vspace{2pt}
      \hallabox{RedRisk}{\faFire\ \ Concentración de Riesgo}{%
        \footnotesize Solo el \textbf{32.7\%} (1.045 personas) en riesgo
        Alto+Muy Alto presenta tasa de rotación real de \textbf{27--30\%}.}
  \end{columns}
\end{frame}

% ─────────────────────────────────────────────────────────────────────────────
\begin{frame}[shrink=5]{Riesgo por Área: Mapa de Calor}
  \begin{columns}[T]
    \column{0.54\textwidth}
      \footnotesize
      \begin{tabular}{@{}lrrr@{}}
        \toprule
        \rowcolor{NavyDeep}
          \color{White}\textbf{Área} &
          \color{White}\textbf{A+MA} &
          \color{White}\textbf{Total} &
          \color{White}\textbf{\%} \\
        \midrule
        \rowcolor{RedRisk!20} Tecnología    & 285 & 800 & 35.6\% \\
        \rowcolor{RedRisk!12} Operaciones   & 236 & 704 & 33.5\% \\
        \rowcolor{Amber!15}   Atención Cli. & 222 & 640 & 34.7\% \\
        \rowcolor{Amber!10}   Riesgo/Cump.  & 190 & 576 & 33.0\% \\
        \rowcolor{Gold!8}     Finanzas      &  61 & 256 & 23.8\% \\
        \rowcolor{GreenOK!8}  Personas      &  24 & 128 & 18.8\% \\
        \rowcolor{GreenOK!5}  Legal         &  22 &  64 & 34.4\% \\
        \rowcolor{GreenOK!5}  Admin.        &   5 &  32 & 15.6\% \\
        \bottomrule
      \end{tabular}

    \column{0.43\textwidth}
      \vspace{4pt}
      \begin{tcolorbox}[
        enhanced, colback=NavyDeep!5, colframe=NavyLight,
        arc=5pt, boxrule=1pt, left=3pt, right=3pt, top=4pt, bottom=4pt
      ]
        \centering
        {\color{NavyDeep}\footnotesize\bfseries Personas con Riesgo Alto/Muy Alto}
        \vspace{3pt}

        \begin{tikzpicture}
          \begin{axis}[
            ybar, width=5.2cm, height=3.8cm,
            bar width=9pt,
            symbolic x coords={Otras,Riesgo,Clientes,Oper.,Tecnol.},
            xtick={Otras,Riesgo,Clientes,Oper.,Tecnol.},
            enlarge x limits=0.15,
            ymin=0, ymax=45,
            ytick={0,10,20,30,40},
            yticklabels={0\%,10\%,20\%,30\%,40\%},
            tick label style={font=\tiny},
            x tick label style={font=\tiny, rotate=30, anchor=east},
            axis lines=left,
            ymajorgrids=true, grid style={dotted,gray!25},
            title={\tiny \% Pers. Alto/Muy Alto Riesgo},
            title style={font=\tiny,color=NavyMid}
          ]
            \addplot[fill=NavyLight!60,draw=none, bar shift=0pt]
              coordinates {(Otras,21.0)};
            \addplot[fill=Gold!80,draw=none, bar shift=0pt]
              coordinates {(Riesgo,33.0)(Clientes,34.7)};
            \addplot[fill=RedRisk!80,draw=none, bar shift=0pt]
              coordinates {(Oper.,33.5)(Tecnol.,35.6)};
          \end{axis}
        \end{tikzpicture}
      \end{tcolorbox}
  \end{columns}
\end{frame}

% ============================================================================
\section[Modelo ML]{Modelo Predictivo ML}
% ============================================================================

\begin{frame}[shrink=5]{Comparación de Modelos de Machine Learning}
  \begin{columns}[T]
    \column{0.52\textwidth}
      \footnotesize
      \begin{tabular}{@{}lrrrr@{}}
        \toprule
        \rowcolor{NavyDeep}
          \color{White}\textbf{Modelo} &
          \color{White}\textbf{Prec.} &
          \color{White}\textbf{AUC} &
          \color{White}\textbf{Sens.} &
          \color{White}\textbf{Espec.} \\
        \midrule
        Reg.\ Logística & 78\% & 0.677 & 1.0\% & 100\% \\
        \rowcolor{GreenOK!20}
        \textbf{Random Forest} \faAward
          & \textbf{78\%} & \textbf{0.659} & \textbf{0.0\%} & \textbf{100\%} \\
        \bottomrule
      \end{tabular}
      \vspace{8pt}

      \begin{tcolorbox}[
        enhanced, colback=GreenOK!8, colframe=GreenOK,
        arc=3pt, boxrule=0.8pt, left=4pt, right=4pt, top=2pt, bottom=2pt
      ]
        {\color{GreenOK}\bfseries\footnotesize \faCheckCircle\ \ Modelo Seleccionado:}\\
        {\footnotesize\textbf{Random Forest} — AUC 0.659.
        Seleccionado por robustez e interpretabilidad en segmentación.}
      \end{tcolorbox}
      \vspace{2pt}
      {\footnotesize\textbf{\color{NavyDeep}Matriz de Confusión}}\\[2pt]
      \footnotesize
      \begin{tabular}{@{}r|cc@{}}
        & \textbf{Pred.\ No} & \textbf{Pred.\ Sí} \\
        \hline
        \textbf{Real No} & \cellcolor{GreenOK!20}499 & \cellcolor{Amber!30}0 \\
        \textbf{Real Sí} & \cellcolor{RedRisk!20}141   & \cellcolor{GreenOK!20}0 \\
      \end{tabular}

    \column{0.44\textwidth}
      \resizebox{\linewidth}{!}{\bfseries\color{NavyDeep}Importancia de Variables (Random Forest)}
      \vspace{3pt}
      \begin{tikzpicture}
        \begin{axis}[
          xbar, width=3.4cm, height=5.4cm,
          bar width=8pt,
          symbolic y coords={Movilidad,Capacitación,Área Crit.,Compromiso,Ausentismo,Desempeño,Clima,Antigüedad},
          ytick={Movilidad,Capacitación,Área Crit.,Compromiso,Ausentismo,Desempeño,Clima,Antigüedad},
          xmin=0, xmax=22,
          enlarge y limits=0.07,
          xtick={0,5,10,15,20},
          tick label style={font=\tiny},
          yticklabel style={font=\tiny, anchor=east},
          axis x line=bottom,
          axis y line=left,
          xmajorgrids=true, grid style={dotted,gray!25},
          xlabel={\tiny Importancia (\%)},
          xlabel style={font=\tiny, yshift=-2pt}
        ]
          \addplot[fill=NavyLight!70,draw=none, bar shift=0pt] coordinates {
            (1.42,Movilidad)(7.16,Capacitación)(12.14,Área Crit.)
            (14.39,Compromiso)(14.44,Ausentismo)
          };
          \addplot[fill=Amber!80,draw=none, bar shift=0pt] coordinates {
            (15.19,Desempeño)(15.48,Clima)
          };
          \addplot[fill=RedRisk!80,draw=none, bar shift=0pt] coordinates {
            (19.78,Antigüedad)
          };
        \end{axis}
      \end{tikzpicture}
  \end{columns}
\end{frame}

% ─────────────────────────────────────────────────────────────────────────────
\begin{frame}[shrink=3]{Curva ROC y Capacidad Predictiva}
  \begin{columns}[T]
    \column{0.54\textwidth}
      \centering
      {\small\bfseries\color{NavyDeep}Curva ROC Comparativa}
      \begin{tikzpicture}
        \begin{axis}[
          width=7cm, height=5cm,
          xmin=0, xmax=1, ymin=0, ymax=1,
          xlabel={\scriptsize Tasa Falsos Positivos (FPR)},
          ylabel={\scriptsize Tasa Verdaderos Positivos (TPR)},
          xlabel style={font=\scriptsize},
          ylabel style={font=\scriptsize},
          tick label style={font=\tiny},
          axis lines=left,
          grid=both, grid style={dotted,gray!25},
          legend style={font=\tiny, at={(0.95,0.15)}, anchor=south east}
        ]
          % Diagonal aleatoria
          \addplot[gray,dashed,thin] coordinates {(0,0)(1,1)};
          % Random Forest (AUC=0.659)
          \addplot[NavyLight,thick,smooth] coordinates {
            (0,0)(0.05,0.20)(0.10,0.35)(0.20,0.52)
            (0.30,0.63)(0.50,0.75)(0.70,0.86)(1.0,1.0)
          };
          % Reg. Logística (AUC=0.677)
          \addplot[Gold,thick,smooth,dashed] coordinates {
            (0,0)(0.05,0.22)(0.10,0.37)(0.20,0.54)
            (0.30,0.65)(0.50,0.77)(0.70,0.87)(1.0,1.0)
          };
          \legend{Azar, RF (AUC=0.66), LogReg (AUC=0.68)}
        \end{axis}
      \end{tikzpicture}

    \column{0.42\textwidth}
      \vspace{3pt}
      \begin{block}{\faInfoCircle\ \ ¿Qué significa AUC 0.68?}
        \small
        El modelo AUC 0.68 discrimina mejor que el azar:
        ordena personas por probabilidad de fuga con
        68\% de acierto al comparar pares rotador/no rotador.
      \end{block}
      \vspace{3pt}
      \hallabox{GreenOK}{\faBullseye\ \ Segmentación de Riesgo}{%
        Con AUC 0.677, el modelo ordena a los 3.200
        colaboradores por nivel de riesgo, priorizando
        intervenciones en el \textbf{32.7\%} en zona
        Alta+Muy Alta.}
  \end{columns}
\end{frame}

% ============================================================================
\section[Plan de Acción]{Plan de Acción}
% ============================================================================

\begin{frame}{Matriz de Priorización: Riesgo × Impacto}
  \begin{columns}[T]
    \column{0.52\textwidth}
      \centering
      \begin{tikzpicture}
        \begin{axis}[
          width=6.5cm, height=5.5cm,
          xmin=0, xmax=110, ymin=0, ymax=110,
          xlabel={\scriptsize Score de Riesgo (0--100)},
          ylabel={\scriptsize Impacto Organizacional (0--100)},
          xlabel style={font=\scriptsize},
          ylabel style={font=\scriptsize},
          tick label style={font=\tiny},
          axis lines=left, grid=both, grid style={dotted,gray!20}
        ]
          % Zonas de color
          \fill[GreenOK,opacity=0.08] (axis cs:0,0) rectangle (axis cs:50,50);
          \fill[Gold,opacity=0.12]    (axis cs:50,0) rectangle (axis cs:75,75);
          \fill[Amber,opacity=0.12]   (axis cs:0,50) rectangle (axis cs:50,100);
          \fill[RedRisk,opacity=0.15] (axis cs:75,75) rectangle (axis cs:100,100);
          % Etiquetas de zona
          \node[font=\tiny,color=GreenOK!70!black] at (axis cs:25,25) {BAJO};
          \node[font=\tiny,color=Amber!70!black]   at (axis cs:25,75) {MEDIO};
          \node[font=\tiny,color=Gold!70!black]    at (axis cs:62,35) {MEDIO-A};
          \node[font=\tiny,color=RedRisk]          at (axis cs:87,87) {\bfseries CRÍTICO};
          % Puntos de ejemplo – tres series con colores fijos
          \addplot[only marks, mark=*, mark size=2.5pt,
            draw=GreenOK, fill=GreenOK!70, opacity=0.9
          ] coordinates {
            (20,30)(15,45)(30,20)(40,35)
          };
          \addplot[only marks, mark=*, mark size=2.5pt,
            draw=Gold, fill=Gold!80, opacity=0.9
          ] coordinates {
            (55,60)(65,40)(70,70)(60,55)
          };
          \addplot[only marks, mark=*, mark size=2.5pt,
            draw=RedRisk, fill=RedRisk!80, opacity=0.9
          ] coordinates {
            (80,85)(90,80)(85,90)(95,75)(75,80)(88,88)
          };
        \end{axis}
      \end{tikzpicture}

    \column{0.44\textwidth}
      \vspace{4pt}
      \small
      \begin{tabular}{@{}lrr@{}}
        \toprule
        \rowcolor{NavyDeep}
          \color{White}\textbf{Zona} &
          \color{White}\textbf{N} &
          \color{White}\textbf{Acción} \\
        \midrule
        \rowcolor{RedRisk!20}\textbf{Muy Alto} &  79 & \textbf{Inmediata} \\
        \rowcolor{Amber!15}  Alto           & 966 & Urgente \\
        \rowcolor{Gold!10}   Medio          & 1.385 & Preventiva \\
        \rowcolor{GreenOK!10}Bajo           & 724 & Monitoreo \\
        \bottomrule
      \end{tabular}
      \vspace{8pt}
      \hallabox{RedRisk}{\faRocket\ \ Acción Prioritaria}{%
        Las \textbf{79 personas en Riesgo Muy Alto}
        (2.5\%) presentan una tasa de rotación real del \textbf{30.4\%}.
        Intervenir aquí primero maximiza
        el ROI de las acciones de retención.}
  \end{columns}
\end{frame}

% ─────────────────────────────────────────────────────────────────────────────
\begin{frame}{Proyección Financiera: 3 Escenarios}
  \begin{columns}[T]
    \column{0.56\textwidth}
      \small
      \begin{tabular}{@{}lrrr@{}}
        \toprule
        \rowcolor{NavyDeep}
          \color{White}\textbf{Concepto} &
          \color{White}\textbf{Año 1} &
          \color{White}\textbf{Año 2} &
          \color{White}\textbf{Año 3} \\
        \midrule
        \multicolumn{4}{@{}l}{\cellcolor{RedRisk!10}
          \textit{Escenario A: Sin Intervención}} \\
        Costo rotación  & \$2.1M & \$2.1M & \$2.1M \\
        \midrule
        \multicolumn{4}{@{}l}{\cellcolor{GreenOK!10}
          \textit{Escenario B: Con Intervención}} \\
        Ahorro Directo   & \$588K & \$756K & \$840K \\
        Costo Intervenc. & (\$263K)& (\$150K)& (\$100K) \\
        \rowcolor{GreenOK!20}
        \textbf{Beneficio Neto}
          & \textbf{\$325K} & \textbf{\$606K} & \textbf{\$740K} \\
        \midrule
        \rowcolor{NavyLight!15}
        \textbf{Acumulado}
          & \textbf{\$325K} & \textbf{\$931K} & \textbf{\$1.67M} \\
        \bottomrule
      \end{tabular}
      \vspace{6pt}
      \[
        \text{ROI Año 1} = \frac{\$325.500}{\$262.500} \times 100 = \mathbf{123.9\%}
      \]

    \column{0.40\textwidth}
      {\small\bfseries\color{NavyDeep}Beneficio Neto Acumulado}
      \begin{tikzpicture}
        \begin{axis}[
          ybar, width=5cm, height=4.5cm,
          bar width=18pt,
          xtick={1,2,3},
          xticklabels={Año 1, Año 2, Año 3},
          ymin=0, ymax=2000,
          ytick={0,500,1000,1500,2000},
          yticklabels={\$0,\$0.5M,\$1M,\$1.5M,\$2M},
          tick label style={font=\tiny},
          axis lines=left,
          ymajorgrids=true, grid style={dotted,gray!30}
        ]
          \addplot[fill=GreenOK!80,draw=none]
            coordinates {(1,325)(2,931)(3,1671)};
        \end{axis}
      \end{tikzpicture}\\
      \hallabox{GreenOK}{\faPiggyBank\ \ USD 1.67M en 3 años}{%
        Por cada \textbf{USD 1} invertido en
        retención se obtienen \textbf{USD 2.24}
        en beneficios netos.}
  \end{columns}
\end{frame}

% ─────────────────────────────────────────────────────────────────────────────
\begin{frame}{Plan de Intervención: 3 Niveles}
  \begin{tikzpicture}[overlay,remember picture]
    \fill[GrayLight,rounded corners=4pt]
      ([xshift=6pt,yshift=-74pt]current page.north west)
      rectangle
      ([xshift=-6pt,yshift=18pt]current page.south east);
  \end{tikzpicture}

  \begin{columns}[T]
    \column{0.32\textwidth}
      \begin{tcolorbox}[
        enhanced, colback=RedRisk!8, colframe=RedRisk,
        arc=5pt, boxrule=1.5pt,
        title={\faExclamationCircle\ \ Nivel 1 — Urgente},
        fonttitle=\small\bfseries,
        attach boxed title to top center={yshift=-2mm},
        boxed title style={colback=RedRisk,colframe=RedRisk}
      ]
        {\scriptsize\bfseries Próximas 2 Semanas}\\[3pt]
        {\scriptsize
        \faCircle\ Reuniones 1-a-1 con 79 personas en riesgo muy alto\\[2pt]
        \faCircle\ Diagnóstico rápido de clima en áreas críticas\\[2pt]
        \faCircle\ Análisis de brecha salarial vs.\ mercado\\[4pt]
        \textbf{Inversión: USD \$12.500}
        }
      \end{tcolorbox}

    \column{0.32\textwidth}
      \begin{tcolorbox}[
        enhanced, colback=Amber!8, colframe=Amber,
        arc=5pt, boxrule=1.5pt,
        title={\faShieldAlt\ \ Nivel 2 — Preventivo},
        fonttitle=\small\bfseries,
        attach boxed title to top center={yshift=-2mm},
        boxed title style={colback=Amber,colframe=Amber}
      ]
        {\scriptsize\bfseries Próximo Mes}\\[3pt]
        {\scriptsize
        \faCircle\ Mentoría estructurada para 546 personas en riesgo alto\\[2pt]
        \faCircle\ Plan de capacitación: meta 3+ cursos/año\\[2pt]
        \faCircle\ Talleres de liderazgo para jefaturas\\[4pt]
        \textbf{Inversión: USD \$205.000}
        }
      \end{tcolorbox}

    \column{0.32\textwidth}
      \begin{tcolorbox}[
        enhanced, colback=GreenOK!8, colframe=GreenOK,
        arc=5pt, boxrule=1.5pt,
        title={\faChartBar\ \ Nivel 3 — Transformación},
        fonttitle=\small\bfseries,
        attach boxed title to top center={yshift=-2mm},
        boxed title style={colback=GreenOK,colframe=GreenOK}
      ]
        {\scriptsize\bfseries 6 Meses}\\[3pt]
        {\scriptsize
        \faCircle\ Dashboard ejecutivo de rotación en tiempo real\\[2pt]
        \faCircle\ Integración de People Analytics en decisiones de HR\\[2pt]
        \faCircle\ Revisión de políticas de carrera y compensación\\[4pt]
        \textbf{Inversión: USD \$45.000}
        }
      \end{tcolorbox}
  \end{columns}
\end{frame}

% ─────────────────────────────────────────────────────────────────────────────
\begin{frame}[shrink=8]{Roadmap de Implementación}
  \centering
  \begin{tikzpicture}[
    node distance=0.4cm,
    fase/.style={
      draw=NavyLight, fill=NavyLight!10, rounded corners=4pt,
      text width=2.65cm, align=center, font=\scriptsize, minimum height=1.1cm,
      inner sep=2pt
    },
    arrow/.style={-{Stealth[length=5pt]}, NavyLight, thick}
  ]
    \node[fase,fill=RedRisk!10,draw=RedRisk] (s1) {
      \textbf{\color{RedRisk}Semana 1--2}\\
      Validación ejecutiva\\
      Aprobación presupuesto\\
      Comunicación liderazgo
    };
    \node[fase,fill=Amber!10,draw=Amber,right=0.15cm of s1] (s2) {
      \textbf{\color{Amber}Semana 2--4}\\
      Equipos de gestión\\
      Capacitar jefaturas\\
      Reuniones 1-a-1
    };
    \node[fase,fill=Gold!10,draw=Gold,right=0.15cm of s2] (m2) {
      \textbf{\color{Gold}Mes 2--3}\\
      Programa mentoría\\
      Encuesta de clima\\
      Plan capacitación
    };
    \node[fase,fill=NavyLight!10,draw=NavyLight,right=0.15cm of m2] (m4) {
      \textbf{\color{NavyLight}Mes 4--6}\\
      Dashboard BI\\
      Revisión políticas\\
      Analytics integrado
    };
    \node[fase,fill=GreenOK!10,draw=GreenOK,right=0.15cm of m4] (m6) {
      \textbf{\color{GreenOK}Mes 6--12}\\
      Medición de impacto\\
      Ajuste de modelos\\
      Reporte C-Suite
    };

    \draw[arrow] (s1) -- (s2);
    \draw[arrow] (s2) -- (m2);
    \draw[arrow] (m2) -- (m4);
    \draw[arrow] (m4) -- (m6);

    % Línea de tiempo
    \draw[NavyDeep!30, thick, dashed]
      ([yshift=-0.5cm]s1.south west) -- ([yshift=-0.5cm]m6.south east);
    \node[font=\tiny,color=GrayMid] at ([yshift=-0.7cm]s1.south)  {Inicio};
    \node[font=\tiny,color=GrayMid] at ([yshift=-0.7cm]m6.south)  {12 meses};

  \end{tikzpicture}

  \vspace{2pt}
  \hallabox{NavyLight}{\faFlag\ \ Hito Clave — Mes 3}{%
    Primera medición de impacto: revisión de tasa de rotación, NPS interno y evolución del
    score de riesgo. Ajuste de intervenciones según resultados.}
\end{frame}

% ============================================================================
\section[Conclusiones]{Conclusiones}
% ============================================================================

\begin{frame}[shrink=25]{Hallazgos Principales}
  \begin{columns}[T]
    \column{0.48\textwidth}
      \hallabox{RedRisk}{\faExclamationTriangle\ \ Crisis en Áreas Críticas}{%
        Rotación 25.2\% vs.\ 3.8\% en no-críticas.
        Diferencial de \textbf{6.6×}. Riesgo operacional inminente
        en Tecnología y Atención Cli.}
      \vspace{1pt}
      \hallabox{Amber}{\faSearch\ \ Predictores Claros}{%
        Antigüedad (19.8\%), Clima (15.5\%) y Ausentismo (14.4\%) lideran según RF.
        La probabilidad promedio de fuga por área (modelo RF)
        permite focalizar intervenciones de forma eficiente.}
      \vspace{1pt}
      \hallabox{NavyLight}{\faRobot\ \ Modelo Confiable}{%
        AUC 0.677 permite segmentación y priorización
        con un 32.7\% en zona de alto riesgo.}

    \column{0.48\textwidth}
      \hallabox{Gold}{\faBullseye\ \ Foco de Alta Eficiencia}{%
        Solo \textbf{79 personas (2.5\%)} en riesgo Muy Alto
        presentan tasa de rotación real del 30.4\%.}
      \vspace{1pt}
      \hallabox{GreenOK}{\faDollarSign\ \ ROI Comprobado}{%
        Intervención estratégica genera \textbf{USD 1.67M} en
        beneficios netos a 3 años. ROI Año 1: 123.9\%.}
      \vspace{1pt}
      \hallabox{Cyan}{\faClock\ \ Urgencia de Acción}{%
        Cada mes de inacción cuesta \textbf{USD 175.750}
        (USD 2.109.000 / 12). La ventana de acción
        preventiva es ahora.}
  \end{columns}
\end{frame}

% ─────────────────────────────────────────────────────────────────────────────
{
\setbeamertemplate{headline}{}
\setbeamertemplate{footline}{}
\begin{frame}[plain]
  \begin{tikzpicture}[overlay,remember picture]
    \fill[NavyDeep] (current page.south west) rectangle (current page.north east);
    \fill[NavyMid,opacity=0.4]
      (current page.north west) --
      ([xshift=10cm]current page.north west) --
      (current page.center) -- cycle;
    \draw[Cyan,line width=1.5pt]
      ([yshift=-0.35\paperheight]current page.north west) --
      ([yshift=-0.35\paperheight]current page.north east);
    \draw[Gold,line width=0.5pt,opacity=0.5]
      ([yshift=-0.355\paperheight]current page.north west) --
      ([yshift=-0.355\paperheight,xshift=0.5\paperwidth]current page.north west);
  \end{tikzpicture}

  \vspace{1.8cm}
  \centering
  {\color{GrayMid}\small\bfseries\MakeUppercase{FinanCorp Chile · Gerencia de Personas}}
  \vspace{0.5cm}

  {\color{White}\fontsize{22}{26}\selectfont\bfseries
    ¿Preguntas?}
  \vspace{0.3cm}

  {\color{Cyan}\fontsize{15}{18}\selectfont
    Gracias por su atención}
  \vspace{0.8cm}

  \begin{tikzpicture}
    \draw[Cyan,line width=1pt] (0,0) -- (5,0);
    \draw[Gold,line width=0.4pt] (0,-3pt) -- (2.5,-3pt);
  \end{tikzpicture}

\end{frame}
}

\end{document}
